\begin{frame}{Motivation: LLMs are bad calculators}
\begin{columns}[T,totalwidth=\textwidth]
\begin{column}{0.54\textwidth}
\begin{block}{Observation behind \PAL{}}
Large language models can often translate a natural-language problem into useful intermediate steps, but may still make arithmetic or logical execution mistakes.
\end{block}
\vspace{0.15cm}
\begin{itemize}
  \item Few-shot prompting shows examples at test time.
  \item Chain-of-thought prompting asks the model to write intermediate reasoning steps.
  \item But in \CoT{}, the same model both \emph{plans} and \emph{executes} the solution.
\end{itemize}
\end{column}
\begin{column}{0.42\textwidth}
\begin{alertblock}{Problem}
Even if the decomposition is plausible, the final answer can be wrong because the LLM performs calculations as text generation rather than exact computation.
\end{alertblock}
\vspace{0.1cm}
\centering
\begin{tikzpicture}[node distance=0.28cm, every node/.style={font=\scriptsize}]
\node[draw,rounded corners,fill=blue!8,minimum width=3.6cm,align=center] (q) {Natural-language\ problem};
\node[draw,rounded corners,fill=blue!8,below=of q,minimum width=3.6cm,align=center] (steps) {Reasoning steps};
\node[draw,rounded corners,fill=red!10,below=of steps,minimum width=3.6cm,align=center] (calc) {LLM executes\ calculation};
\node[draw,rounded corners,fill=red!10,below=of calc,minimum width=3.6cm,align=center] (ans) {Final answer};
\draw[-{Latex}] (q) -- (steps);
\draw[-{Latex}] (steps) -- (calc);
\draw[-{Latex}] (calc) -- (ans);
\end{tikzpicture}
\end{column}
\end{columns}
% \source{PAL paper, Abstract and Introduction.}
\end{frame}
